% Source: https://artofproblemsolving.com/community/c3h2192694_54152674858
% Stable identity: aops-topic-2192694-q01
\textbf{(a)} Place the balls in the order $1,2,\ldots,2020$. Ball $1$ can go into any of the $3$ boxes. For $k\ge2$, the only condition involving ball $k$ and the balls before it is that ball $k$ is not in the same box as ball $k-1$; so ball $k$ has exactly $2$ choices, whatever happened before. Every valid distribution arises exactly once, so the number of ways is
\[\boxed{3\cdot2^{2019}}.\]

\textbf{(b)} From the distributions in (a) remove those that leave a box empty. All the balls cannot go into one box (balls $1$ and $2$ would be together). If exactly two boxes are used, the balls must alternate between them, since consecutive balls are in different boxes: once the box of ball $1$ is chosen, everything is determined. There are $3$ choices of the pair of boxes and $2$ choices of the box of ball $1$, so $6$ such distributions. Hence the number of ways is
\[\boxed{3\cdot2^{2019}-6}.\]
